% Source coverage: physical PDF pp. 248--257 (printed pp. 225--234), beginning at the Section 19.7 heading and ending immediately before the Section 19.8 heading.
% Numbered equations: (19.7.1)--(19.7.32). Citations: [33]--[41], including [38a] and [38b]. Footnotes: 2.
% Uncertainty: none.

\subsection{\texorpdfstring{Effective Field Theories: $SU(3)\cdot SU(3)$}{Effective Field Theories: SU(3) . SU(3)}}
\label{sec:19.7}

The approximate $SU(2)$ isospin symmetry of nuclear physics was extended in the early 1960s by Gell-Mann~\hyperref[ch19-ref-33]{[33]} and Ne'eman~\hyperref[ch19-ref-34]{[34]} to an even less exact $SU(3)$ symmetry, which grouped the known baryons and mesons in various irreducible representations: an octet of $\frac{1}{2}^+$ baryons $p,n,\Lambda^0,\Sigma^{\pm,0},\allowbreak\Xi^{-,0}$, an octet of $0^-$ mesons $K^{+,0},\pi^{\pm,0},\eta^0,\allowbreak\overline K^{-,0}$, an octet of $1^-$ mesons $K^{*+,0},\rho^{\pm,0},\omega,\allowbreak\overline K^{*-,0}$, and a decuplet of $\frac{3}{2}^+$ baryons $\Delta^{++,+,0,-},\Sigma^{*+,0,-},\allowbreak\Xi^{*0,-},\Omega^-$. (The $\eta$ and $\Omega^-$ were not discovered until later.) After the successes of the chiral $SU(2)\cdot SU(2)$ symmetry in the mid-1960s, it became natural to suppose that the strong interactions also respect an approximate $SU(3)\cdot SU(3)$ symmetry, which like $SU(2)\cdot SU(2)$ is spontaneously broken to its diagonal subgroup, the $SU(3)$ of Gell-Mann and Ne'eman. Then after the advent of quantum chromodynamics it became clear that this symmetry arises because there are not two but three fairly light quarks: the $u$ and $d$, and a third quark $s$ which like $d$ has charge $-1/3$. In this case the $SU(3)\cdot SU(3)$ symmetry consists of independent $SU(3)$ transformations (analogous to Eq.~\eqref{eq:19.4.2}) on the left- and right-handed parts of the $u,d,s$ quark fields:
\begin{equation}
\begin{pmatrix}u\\d\\s\end{pmatrix}
\longrightarrow
\exp\left[i\sum_a\left(\theta_a^V\lambda_a+
\theta_a^A\lambda_a\gamma_5\right)\right]
\begin{pmatrix}u\\d\\s\end{pmatrix}.
\tag{19.7.1}\label{eq:19.7.1}
\end{equation}
Here the $\lambda_a$ are a complete set of traceless Hermitian matrices:
\begin{equation}
\begin{gathered}
\lambda_1=\begin{pmatrix}0&1&0\\1&0&0\\0&0&0\end{pmatrix},
\qquad
\lambda_2=\begin{pmatrix}0&-i&0\\i&0&0\\0&0&0\end{pmatrix},
\qquad
\lambda_3=\begin{pmatrix}1&0&0\\0&-1&0\\0&0&0\end{pmatrix},
\\[6pt]
\lambda_4=\begin{pmatrix}0&0&1\\0&0&0\\1&0&0\end{pmatrix},
\qquad
\lambda_5=\begin{pmatrix}0&0&-i\\0&0&0\\i&0&0\end{pmatrix},
\qquad
\lambda_6=\begin{pmatrix}0&0&0\\0&0&1\\0&1&0\end{pmatrix},
\\[6pt]
\lambda_7=\begin{pmatrix}0&0&0\\0&0&-i\\0&i&0\end{pmatrix},
\qquad
\lambda_8=\frac{1}{\sqrt3}
\begin{pmatrix}1&0&0\\0&1&0\\0&0&-2\end{pmatrix}.
\end{gathered}
\tag{19.7.2}\label{eq:19.7.2}
\end{equation}
They are normalized so that $\operatorname{Tr}(\lambda_a\lambda_b)=2\delta_{ab}$. The generators of $SU(3)\cdot SU(3)$ are thus represented on the quark fields by the generators $t_a=\lambda_a$ of the unbroken $SU(3)$ symmetry and the broken symmetry generators $x_a=\lambda_a\gamma_5$.

To define the Goldstone boson fields and work out their transformation properties, we note that in the representation furnished by the quark fields, any $SU(3)\cdot SU(3)$ transformation may be written as $\exp(-i\gamma_5\sum_a\xi_a\lambda_a)$ times a transformation $\exp(i\sum_a\theta_a\lambda_a)$ belonging to the unbroken $SU(3)$ subgroup of $SU(3)\cdot SU(3)$. (The minus sign accompanying $\xi$ is inserted for convenience in later comparing components of this field with the pion fields introduced in Section 19.5.) Hence in this representation, each right coset of $SU(3)\cdot SU(3)/SU(3)$ is represented by the matrix of the form $\gamma(\xi)=\exp(-i\gamma_5\sum_a\xi_a\lambda_a)$ that it contains, and aside from normalization, the parameters $\xi_a$ of these right cosets may be taken as our Goldstone boson fields. According to Eq.~\eqref{eq:19.6.18}, the transformation rule of these fields is dictated by
\begin{equation}
\begin{split}
&\exp\left[i\sum_a\left(\theta_a^V\lambda_a+
\theta_a^A\lambda_a\gamma_5\right)\right]
\exp\left[-i\gamma_5\sum_a\xi_a(x)\lambda_a\right]
\\
&\qquad
=\exp\left[-i\gamma_5\sum_a\xi'_a(x)\lambda_a\right]
\exp\left[i\sum_a\theta_a(x)\lambda_a\right],
\end{split}
\tag{19.7.3}\label{eq:19.7.3}
\end{equation}
with $\theta_a(x)$ some function of $\theta^V,\theta^A$, and $\xi(x)$. Also, according to Eq.~\eqref{eq:19.6.3}, the Goldstone-free quark fields $\widetilde q(x)$ are here defined by
\begin{equation}
q(x)\equiv
\begin{pmatrix}u(x)\\d(x)\\s(x)\end{pmatrix}
=\exp\left[-i\gamma_5\sum_a\xi_a(x)\lambda_a\right]\widetilde q(x),
\tag{19.7.4}\label{eq:19.7.4}
\end{equation}
and have the transformation rule given by Eq.~\eqref{eq:19.6.24}:
\begin{equation}
\widetilde q'(x)
=\exp\left[i\sum_a\theta_a(x)\lambda_a\right]\widetilde q(x).
\tag{19.7.5}\label{eq:19.7.5}
\end{equation}
We could proceed to introduce covariant derivatives \eqref{eq:19.6.15} and \eqref{eq:19.6.30}, and use them to construct chiral-invariant Lagrangian densities, but for chiral symmetries there is a simpler approach available.

Note that the parts of Eq.~\eqref{eq:19.7.3} that are proportional to $(1+\gamma_5)$ and $(1-\gamma_5)$ read
\begin{equation}
\begin{split}
&\exp\left(i\sum_a\theta_a^L\lambda_a\right)
\exp\left(-i\sum_a\xi_a(x)\lambda_a\right)
\\
&\qquad
=\exp\left(-i\sum_a\xi'_a(x)\lambda_a\right)
\exp\left(i\sum_a\theta_a(x)\lambda_a\right)
\end{split}
\tag{19.7.6}\label{eq:19.7.6}
\end{equation}
and
\begin{equation}
\begin{split}
&\exp\left(i\sum_a\theta_a^R\lambda_a\right)
\exp\left(i\sum_a\xi_a(x)\lambda_a\right)
\\
&\qquad
=\exp\left(i\sum_a\xi'_a(x)\lambda_a\right)
\exp\left(i\sum_a\theta_a(x)\lambda_a\right),
\end{split}
\tag{19.7.7}\label{eq:19.7.7}
\end{equation}
where $\theta^L=\theta^V+\theta^A$ and $\theta^R=\theta^V-\theta^A$. By multiplying Eq.~\eqref{eq:19.7.7} on the right by the inverse of Eq.~\eqref{eq:19.7.6}, we find the simple transformation rule
\begin{equation}
U'(x)
=\exp\left(i\sum_a\lambda_a\theta_a^R\right)
U(x)
\exp\left(-i\sum_a\lambda_a\theta_a^L\right),
\tag{19.7.8}\label{eq:19.7.8}
\end{equation}
where $U(x)$ is the unitary unimodular matrix
\begin{equation}
U(x)=\exp\left(2i\sum_a\xi_a(x)\lambda_a\right).
\tag{19.7.9}\label{eq:19.7.9}
\end{equation}
In other words, $U(x)$ transforms according to the $(\overline{\mathbf 3},\mathbf 3)$ representation of $SU(3)\cdot SU(3)$. This is still a non-linear realization of $SU(3)\cdot SU(3)$ because the components of $U(x)$ are not independent; they are subject to the nonlinear constraints $U^\dagger U=\mathbf 1$ and $\det U=1$.

The unique $(SU(3)\cdot SU(3))$-invariant term in the Goldstone boson Lagrangian that is of second order in spacetime derivatives is
\begin{equation}
\mathcal L_{\text{2 deriv}}
=-\frac{1}{16}F^2\operatorname{Tr}
\left\{\partial_\mu U\,\partial^\mu U^\dagger\right\},
\tag{19.7.10}\label{eq:19.7.10}
\end{equation}
with $F^2$ a positive constant to be determined. We can express the $\xi_a$ in terms of conventionally normalized pseudoscalar meson fields by writing
\begin{equation}
\sum_a\lambda_a\xi_a
=\frac{\sqrt2}{F}
\begin{pmatrix}
\dfrac{\pi^0}{\sqrt2}+\dfrac{\eta^0}{\sqrt6}&\pi^+&K^+\\[5pt]
\pi^-&-\dfrac{\pi^0}{\sqrt2}+\dfrac{\eta^0}{\sqrt6}&K^0\\[5pt]
\overline K^-&\overline K^0&-\sqrt{\dfrac23}\eta^0
\end{pmatrix}
\equiv\frac{\sqrt2}{F}B,
\tag{19.7.11}\label{eq:19.7.11}
\end{equation}
so that the kinetic term in Eq.~\eqref{eq:19.7.10} is of the usual form
\[
\mathcal L_{\text{kin}}
=-\frac12\partial_\mu\pi^0\partial^\mu\pi^0
-\partial_\mu\pi^+\partial^\mu\pi^-
-\partial_\mu K^+\partial^\mu\overline K^-
-\partial_\mu K^0\partial^\mu\overline K^0
-\frac12\partial_\mu\eta^0\partial^\mu\eta^0.
\]
To determine the constant $F$, we note by comparing Eq.~\eqref{eq:19.7.4} with Eq.~\eqref{eq:19.5.42}, that for $\xi$ infinitesimal, the components $\xi_1$, $\xi_2$, and $\xi_3$ are the same as the Goldstone boson fields $\zeta_1$, $\zeta_2$, and $\zeta_3$ introduced in Section 19.5. Then, comparing \eqref{eq:19.7.11} with \eqref{eq:19.5.17} and \eqref{eq:19.5.32}, we see that the constant $F$ defined by Eq.~\eqref{eq:19.7.10} is the same as the $F_\pi=184\ \mathrm{MeV}$ introduced in Section 19.4.

The $SU(3)\cdot SU(3)$ symmetry of quantum chromodynamics is broken by the quark mass term, which in terms of the quark field $\widetilde q(x)$ defined by Eq.~\eqref{eq:19.7.4} is
\begin{equation}
\mathcal L_{\text{mass}}
=-\overline q M_q q
=-\overline{\widetilde q}\,
\exp\left(-i\sqrt2\gamma_5B/F_\pi\right)
M_q
\exp\left(-i\sqrt2\gamma_5B/F_\pi\right)
\widetilde q,
\tag{19.7.12}\label{eq:19.7.12}
\end{equation}
where
\begin{equation}
M_q=
\begin{pmatrix}
m_u&0&0\\
0&m_d&0\\
0&0&m_s
\end{pmatrix}.
\tag{19.7.13}\label{eq:19.7.13}
\end{equation}
Eq.~\eqref{eq:19.7.12} contains a purely bosonic part, obtained by replacing the quark bilinear with its vacuum expectation value, given by the unbroken $SU(3)$ and parity symmetries as\footnote{The fact that the conservation laws of parity, charge, and strangeness respected by the vacuum are the same as those respected by the quark mass matrix is a result of the vacuum alignment condition discussed in Section 19.3. Likewise, if the quark fields are defined so that the quark masses are all positive, then with $v>0$ the minus sign in Eq.~\eqref{eq:19.7.3} is required by the condition that the true vacuum be at a minimum rather than a maximum of the vacuum energy for vacuum states rotated by $SU(3)\cdot SU(3)$ transformations, for as we shall see, it is this sign that will yield positive masses for the pseudo-Goldstone boson octet.}
\begin{equation}
\left\langle\overline{\widetilde q}_n\gamma_5\widetilde q_m\right\rangle_0=0,
\qquad
\left\langle\overline{\widetilde q}_n\widetilde q_m\right\rangle_0
=-v\delta_{nm}.
\tag{19.7.14}\label{eq:19.7.14}
\end{equation}
The Goldstone boson mass term in the Lagrangian is then
\begin{equation}
\begin{split}
-\frac{v}{F_\pi^2}\operatorname{Tr}\{B,\{B,M_q\}\}
=-\frac{v}{F_\pi^2}\bigg[
&4m_u\left(\frac{\pi^0}{\sqrt2}+\frac{\eta^0}{\sqrt6}\right)^2
+4(m_u+m_d)\pi^+\pi^-
\\
&+4(m_u+m_s)K^+\overline K^-
+4m_d\left(-\frac{\pi^0}{\sqrt2}+\frac{\eta^0}{\sqrt6}\right)^2
\\
&+4(m_d+m_s)K^0\overline K^0
+\frac83m_s(\eta^0)^2
\bigg].
\end{split}
\tag{19.7.15}\label{eq:19.7.15}
\end{equation}
From this, we can read off that~\hyperref[ch19-ref-35]{[35]}
\begin{equation}
\begin{aligned}
m_{\pi^+}^2=m_{\pi^0}^2
&=\frac{4v}{F_\pi^2}(m_u+m_d),\\
m_{K^+}^2
&=\frac{4v}{F_\pi^2}(m_u+m_s),\\
m_{K^0}^2
&=\frac{4v}{F_\pi^2}(m_d+m_s),\\
m_{\eta^0}^2
&=\frac{4v}{F_\pi^2}
\left(\frac{4m_s+m_d+m_u}{3}\right),
\end{aligned}
\tag{19.7.16}\label{eq:19.7.16}
\end{equation}
and there is also a term that mixes the $\pi^0$ and $\eta^0$:
\begin{equation}
m_{\pi\eta}^2
=\frac{4v}{\sqrt3F_\pi^2}(m_u-m_d).
\tag{19.7.17}\label{eq:19.7.17}
\end{equation}

The same results can also be obtained by operator methods. Noether's theorem allows us to construct generators $T_a$ and $X_a$, for which
\begin{equation}
[T_a,q]=-\lambda_aq,
\qquad
[X_a,q]=-\gamma_5\lambda_aq.
\tag{19.7.18}\label{eq:19.7.18}
\end{equation}
We can write the Goldstone boson fields in a real basis as $\pi_a$, with $\pi^\pm=(\pi_1\pm i\pi_2)/\sqrt2$, $\pi^0=\pi_3$, $K^\pm=(\pi_4\pm i\pi_5)/\sqrt2$, $K^0=(\pi_6+i\pi_7)/\sqrt2$, $\overline K^0=(\pi_6-i\pi_7)/\sqrt2$, and $\eta^0=\pi_8$. The $SU(3)$ symmetry that survives the spontaneous symmetry breaking tells us that the $F$-matrix for these bosons takes the form $F_{ab}=F_\pi\delta_{ab}$. The $SU(3)\cdot SU(3)$ symmetry is also intrinsically broken by the mass term in the Hamiltonian $H_1=m_u\overline uu+m_d\overline dd+m_s\overline ss=\overline qM_qq$, where $M_q$ is the quark mass matrix \eqref{eq:19.7.13}. The mass matrix of the pseudo-Goldstone bosons is given here by \eqref{eq:19.3.20} as
\begin{equation}
M_{ab}^2
=-F_\pi^{-2}
\left\langle\overline q\{\lambda_a,\{\lambda_b,M_q\}\}q\right\rangle_0.
\tag{19.7.19}\label{eq:19.7.19}
\end{equation}
Since this is already of first order in quark masses, to this order we may use the unbroken $SU(3)$ relations $\langle\overline uu\rangle_0=\langle\overline dd\rangle_0=\langle\overline ss\rangle_0\equiv-v$, and find the results \eqref{eq:19.7.16} and \eqref{eq:19.7.17}, as before.

The masses $m_{K^+}=493.65\ \mathrm{MeV}$ and $m_{K^0}=497.7\ \mathrm{MeV}$ of the kaon doublet are quite close, and considerably larger than the pion mass, so we can see from \eqref{eq:19.7.16} that the $m_u$ and $m_d$ must be considerably smaller than $m_s$. In calculating quantities that are sensitive to $m_u$ and $m_d$, like the kaon mass difference or the pion masses, we should also take into account another small correction, the effect of electromagnetism. The electromagnetic current is $J^\mu=ie\overline q\gamma^\mu Qq$, where $Q$ is a diagonal matrix with elements $2/3,-1/3$, and $-1/3$. Its commutators with the $SU(3)\cdot SU(3)$ generators are
\[
[T_a,J^\mu]
=-\frac{ie}{2}\overline q\gamma^\mu[Q,\lambda_a]q,
\qquad
[X_a,J^\mu]
=-\frac{ie}{2}\overline q\gamma^\mu\gamma_5[Q,\lambda_a]q.
\]
We see that $J^\mu$ commutes with $X_3,X_6,X_7$, and $X_8$ as well as $T_3,T_6,T_7$, and $T_8$, so the electromagnetic part of the Hamiltonian is invariant under the $SU(2)\cdot SU(2)\cdot U(1)\cdot U(1)$ subgroup with these generators. Therefore in the limit of zero quark masses, \emph{electromagnetic effects give no mass to the neutral pseudo-Goldstone bosons}~\hyperref[ch19-ref-36]{[36]} $\pi^0,K^0,\overline K^0$, and $\eta^0$, the Goldstone bosons associated with the spontaneous breakdown of the symmetries with generators $X_3,X_6,X_7$, and $X_8$. Also, in the zero-quark-mass limit there is an unbroken $SU(2)$ symmetry generated by $T_6,T_7$, and $\sqrt3T_8-T_3$, under which the $K^+$ and $\pi^+$ transform as a doublet, so that for zero quark masses \emph{the electromagnetic corrections to the $K^+$ and $\pi^+$ masses are equal}~\hyperref[ch19-ref-36]{[36]}. Since the quark mass terms and electromagnetic corrections are all small, it is reasonable to treat these effects as additive corrections to the effective quark Hamiltonian. We see then that with electromagnetism taken into account, the mass formulas \eqref{eq:19.7.16} should be corrected to read~\hyperref[ch19-ref-37]{[37]}
\begin{equation}
\begin{aligned}
m_{K^\pm}^2
&=\frac{4v(m_u+m_s)}{F_\pi^2}+\Delta,\\
m_{K^0}^2=m_{\overline K^0}^2
&=\frac{4v(m_d+m_s)}{F_\pi^2},\\
m_{\pi^\pm}^2
&=\frac{4v(m_d+m_u)}{F_\pi^2}+\Delta,\\
m_{\pi^0}^2
&=\frac{4v(m_d+m_u)}{F_\pi^2},\\
m_\eta^2
&=\frac{4v(m_u+m_d+4m_s)}{3F_\pi^2},
\end{aligned}
\tag{19.7.20}\label{eq:19.7.20}
\end{equation}
where $\Delta$ is the common electromagnetic correction to the $K^+$ and $\pi^+$ squared masses.

These formulas impose one linear relation on the five pseudo-Goldstone masses, a version of the Gell-Mann--Okubo relation:\footnote{This relation was originally derived~\hyperref[ch19-ref-38]{[38]} on the basis of the approximate $SU(3)$ symmetry of Gell-Mann and Ne'eman generated by the $T_a$, ignoring mass differences within isotopic multiplets. From this derivation, one cannot tell whether this relation should be applied to the pseudoscalar meson masses themselves, or as here to their squares. Applying this relation to the pseudoscalar masses themselves would not in fact work very well; it gives an $\eta$ mass of $613\ \mathrm{MeV}$. The fact that the Gell-Mann--Okubo relation works well for the squares of the meson masses, but not for their first powers, supports the interpretation of these particles as Goldstone bosons of a spontaneously broken approximate $SU(3)\cdot SU(3)$ symmetry. Gell-Mann and Okubo also used the approximate $SU(3)$ symmetry generated by the $T_a$ to derive relations among the masses of other particles, such as (ignoring isospin violating effects) the relation $2m_N+2m_\Xi=3m_\Lambda+m_\Sigma$ among the masses of the lightest baryon octet. For such multiplets, the average mass is so much larger than the mass differences within the multiplet that it makes little difference whether one applies the relation to the masses or their squares.}
\begin{equation}
3m_\eta^2+2m_{\pi^+}^2-m_{\pi^0}^2
=2m_{K^+}^2+2m_{K^0}^2.
\tag{19.7.21}\label{eq:19.7.21}
\end{equation}
Taking the kaon and pion masses from experiment, this relation yields an $\eta$ mass of $566\ \mathrm{MeV}$, in comparison with the experimental value of $547\ \mathrm{MeV}$. The small discrepancy is usually attributed to the mixing of the $\eta$ state with a heavier pseudoscalar particle, the $\eta'$ at $958\ \mathrm{MeV}$.

From the mass formulas \eqref{eq:19.7.20} we may derive formulas for the quark mass ratios in terms of the pion and kaon masses~\hyperref[ch19-ref-37]{[37]}:
\begin{equation}
\frac{m_d}{m_s}
=\frac{m_{K^0}^2+m_{\pi^+}^2-m_{K^+}^2}
{m_{K^0}^2+m_{K^+}^2-m_{\pi^+}^2},
\qquad
\frac{m_u}{m_s}
=\frac{2m_{\pi^0}^2-m_{K^0}^2-m_{\pi^+}^2+m_{K^+}^2}
{m_{K^0}^2+m_{K^+}^2-m_{\pi^+}^2}.
\tag{19.7.22}\label{eq:19.7.22}
\end{equation}
Using the mass values $m_{\pi^+}=139.57\ \mathrm{MeV}$, $m_{\pi^0}=134.974\ \mathrm{MeV}$, $m_{K^+}=493.65\ \mathrm{MeV}$, and $m_{K^0}=497.7\ \mathrm{MeV}$ gives the ratios $m_d/m_s=0.050$ and $m_u/m_s=0.027$. Thus the ratio of the $d$ and $u$ quark masses is closer to 2 than 1. (A 1996 calculation~\hyperref[ch19-ref-38a]{[38a]} using various other pieces of information as well as pseudoscalar meson masses has given values $m_d/m_s=0.053\pm0.002$ and $m_u/m_s=0.029\pm0.003$.)

None of this gives definite values for the individual quark masses. Indeed, these masses are not well defined until we define a renormalization prescription for the quark bilinears. Often this prescription is fashioned so that $m_s$ equals the mass difference between isomultiplets differing by one unit of strangeness in some $SU(3)$ multiplet. For instance, the lightest vector meson octet consists of an isotopic doublet $K^*$ at $892\ \mathrm{MeV}$, interpreted as a bound state of an $s$ antiquark and a $u$ or $d$; its antidoublet, consisting of an $s$ and a $\overline u$ or $\overline d$; and a $T=1$ $\rho$ at $770\ \mathrm{MeV}$ and a $T=0$ $\omega$ at $783\ \mathrm{MeV}$, both interpreted as bound states of a $u$ or $d$ and a $\overline u$ or $\overline d$. If we wish to attribute the difference between the $K^*$ mass and the average of the $\rho$ and $\omega$ masses to the relatively large $s$ quark mass, then we must renormalize the quark bilinear so that
\[
m_s-\frac12(m_u+m_d)
=m_{K^*}-\frac12(m_\rho+m_\omega)
=120\ \mathrm{MeV},
\]
yielding $m_s=125\ \mathrm{MeV}$. With the above quark mass ratios, this gives $m_d=6.0\ \mathrm{MeV}$ and $m_u=3.3\ \mathrm{MeV}$. However, these estimates of mass values are much less reliable than the values of the mass ratios given in Eq.~\eqref{eq:19.7.22}. Often $m_s$ is estimated~\hyperref[ch19-ref-38b]{[38b]} as $180\ \mathrm{MeV}$ rather than $125\ \mathrm{MeV}$.

The mass term \eqref{eq:19.7.12} includes meson-meson interactions. Using Eqs.~\eqref{eq:19.7.14}, \eqref{eq:19.7.11}, and \eqref{eq:19.7.9}, we may write the purely bosonic part of this term in the Lagrangian as
\begin{equation}
\mathcal L_{\text{mass,bosonic}}
=\frac12v\operatorname{Tr}\{M_q(U^\dagger+U)\}.
\tag{19.7.23}\label{eq:19.7.23}
\end{equation}
The form of this term could have been deduced from general symmetry considerations, which also allow us to find the allowed terms of higher order in $M_q$. Suppose we invent an external $3\cdot3$ field $\chi$, and replace the mass term \eqref{eq:19.7.12} in the underlying quantum chromodynamics Lagrangian with the $\chi$-quark coupling term
\begin{equation}
\mathcal L_\chi
=-\overline q\left[
\frac12(1+\gamma_5)\chi
+\frac12(1-\gamma_5)\chi^\dagger
\right]q,
\tag{19.7.24}\label{eq:19.7.24}
\end{equation}
which becomes the same as Eq.~\eqref{eq:19.7.12} when we make the replacements $\chi=\chi^\dagger=M_q$. The point of this procedure is that the Lagrangian \eqref{eq:19.7.24} becomes formally invariant under $SU(3)\cdot SU(3)$ if we give $\chi$ the formal transformation rule
\begin{equation}
\chi\longrightarrow
\exp\left(i\sum_a\lambda_a\theta_a^R\right)
\chi
\exp\left(-i\sum_a\lambda_a\theta_a^L\right).
\tag{19.7.25}\label{eq:19.7.25}
\end{equation}
Thus we can work out the allowed bosonic terms involving quark masses by writing the most general $SU(3)\cdot SU(3)$ Lagrangian (up to some given order in derivatives and $M_q$) involving $U$ and $\chi$, requiring also invariance under the parity transformation
\begin{equation}
U(\boldsymbol x,t)\longleftrightarrow U^\dagger(-\boldsymbol x,t),
\qquad
\chi\longleftrightarrow\chi^\dagger,
\tag{19.7.26}\label{eq:19.7.26}
\end{equation}
and then making the replacements
\begin{equation}
\chi=\chi^\dagger=M_q.
\tag{19.7.27}\label{eq:19.7.27}
\end{equation}
For instance, the interaction $\operatorname{Tr}(U^\dagger\chi+U\chi^\dagger)$ is invariant under $SU(3)\cdot SU(3)$ and parity, and becomes of the same form as Eq.~\eqref{eq:19.7.23} when we give $\chi$ the values \eqref{eq:19.7.27}.

Using this technique, Gasser and Leutwyler~\hyperref[ch19-ref-39]{[39]} have given the complete effective Lagrangian for the pseudoscalar octet of fourth order in momenta or meson masses (with quark masses counted as being of second order in meson masses) as
\begin{equation}
\begin{aligned}
\mathcal L_4={}&
L_1\left[\operatorname{Tr}
\{\partial_\mu U^\dagger\partial^\mu U\}\right]^2
\\
&+L_2\operatorname{Tr}
\{\partial_\mu U\partial_\nu U^\dagger\}
\operatorname{Tr}
\{\partial^\mu U\partial^\nu U^\dagger\}
\\
&+L_3\operatorname{Tr}
\{\partial_\mu U\partial^\mu U^\dagger
\partial_\nu U\partial^\nu U^\dagger\}
\\
&+L_4\operatorname{Tr}
\{\partial_\mu U\partial^\mu U^\dagger\}
\operatorname{Tr}\{M_q(U+U^\dagger)\}
\\
&+L_5\operatorname{Tr}
\{\partial_\mu U\partial^\mu U^\dagger
(M_qU+U^\dagger M_q)\}
\\
&+L_6\left[\operatorname{Tr}
\{M_q(U+U^\dagger)\}\right]^2
\\
&+L_7\left[\operatorname{Tr}
\{(U^\dagger-U)M_q\}\right]^2
\\
&+L_8\operatorname{Tr}
\{(UM_q)^2+(U^\dagger M_q)^2\},
\end{aligned}
\tag{19.7.28}\label{eq:19.7.28}
\end{equation}
where $L_1,\ldots,L_8$ are constants to be determined by comparison with experiment. The complete effective Lagrangian up to fourth order in meson masses and momenta is
\begin{equation}
\mathcal L_{\text{eff}}=\mathcal L_2+\mathcal L_4,
\tag{19.7.29}\label{eq:19.7.29}
\end{equation}
where $\mathcal L_2$ is the sum of the terms \eqref{eq:19.7.10} and \eqref{eq:19.7.23}:
\begin{equation}
\mathcal L_2
=-\frac{1}{16}F^2\operatorname{Tr}
\{\partial_\mu U\partial^\mu U^\dagger\}
+\frac12v\operatorname{Tr}\{M_q(U^\dagger+U)\}.
\tag{19.7.30}\label{eq:19.7.30}
\end{equation}
Electroweak interactions may be included by replacing the derivatives $\partial_\mu$ with suitable gauge-covariant derivatives $D_\mu$, and adding a few additional terms to $\mathcal L_{\text{eff}}$.

Following the same power-counting arguments as in Sections 19.5 and 19.6, to calculate $S$-matrix elements to fourth order in meson masses and momenta we must include both tree graphs to all orders in $\mathcal L_2$ and up to first order in $\mathcal L_4$, and also one-loop graphs constructed from $\mathcal L_2$ alone. Gasser and Leutwyler and others have carried out a comprehensive study of meson dynamics (and associated electroweak interactions) using this effective Lagrangian~\hyperref[ch19-ref-40]{[40]}.

\begin{center}
$* \qquad * \qquad *$
\end{center}

The quark mass term \eqref{eq:19.7.12} naturally affects other $SU(3)$ multiplets. For any multiplet other than the pseudoscalar octet this can be treated as a first-order perturbation, and therefore gives a shift in the mass matrix of a generic multiplet $\ket{i}$ equal to
\begin{equation}
\delta m_{ij}
=\bra{i}\overline qM_qq\ket{j},
\tag{19.7.31}\label{eq:19.7.31}
\end{equation}
and $SU(3)$ may be used~\hyperref[ch19-ref-41]{[41]} to relate the various matrix elements $\bra{i}\overline q_rq_s\ket{j}$. In this way, it is straightforward to show that
\begin{equation}
\frac{\delta m_p-\delta m_n}{m_u-m_d}
=\frac{m_\Xi-m_\Sigma}{m_s}.
\tag{19.7.32}\label{eq:19.7.32}
\end{equation}
Using this together with Eq.~\eqref{eq:19.7.22} gives the quark mass contribution to the nucleon mass splitting $\delta m_p-\delta m_n\simeq-2.5\ \mathrm{MeV}$. This result may be used in Eqs.~\eqref{eq:19.5.65} and \eqref{eq:19.5.66} to calculate the leading violations of isospin symmetry in low energy pion-nucleon interactions. Of course, as defined here $\delta m_p-\delta m_n$ is not the full proton-neutron mass difference, which receives an important contribution also from photon emission and absorption. Because the neutron is electrically neutral, this electromagnetic term is almost certainly positive, in agreement with the fact that the observed proton-neutron mass difference is $-1.3\ \mathrm{MeV}$, leaving $+1.2\ \mathrm{MeV}$ to be accounted for by electromagnetism. Unfortunately, accurate calculation of this electromagnetic mass difference has proved difficult.
